\chapter{Kesimpulan dan Saran}

Bab ini menutup naskah dengan menyajikan kesimpulan penelitian dan saran untuk 
pengembangan lanjutan. Subbab \ref{sec:kesimpulan} merangkum jawaban atas tiga rumusan 
masalah yang telah ditetapkan pada Bab I, disusun berurutan mengikuti alur dari 
ketersediaan data hingga penerapan sistem. Subbab \ref{sec:saran} menguraikan arah 
pengembangan lanjutan yang diturunkan dari batasan penelitian dan keterbatasan yang 
teridentifikasi sepanjang proses penelitian.

\section{Kesimpulan}
\label{sec:kesimpulan}

Penelitian ini menjawab tiga rumusan masalah yang ditetapkan pada Bab I melalui 
pembentukan dataset, perbandingan dan penyetelan model deteksi, serta implementasi sistem 
pemantauan berbasis web. Ketiga jawaban tersebut diuraikan secara berurutan pada subbab 
ini, mengikuti urutan rumusan masalah pada Bab I agar keterkaitan antara permasalahan dan 
temuan dapat ditelusuri secara langsung. Subbab ini diakhiri dengan satu paragraf yang 
merangkum keterkaitan antara ketiga jawaban tersebut sebagai satu rangkaian yang saling 
menopang.

Pertama, dataset deteksi objek untuk sembilan kelas khas kawasan Malioboro berhasil 
dibangun melalui anotasi semi-otomatis berbantuan model \textit{pretrained} dengan koreksi 
manual, menjawab rumusan masalah pertama yang menyatakan dataset semacam ini belum 
tersedia pada sumber data publik seperti COCO. Dataset akhir memuat 1.581 citra dengan 
sembilan kelas, yaitu andong, bajaj, becak, sepeda, bus, mobil, motor, orang, dan truk, 
dengan distribusi yang sangat tidak seimbang antara kelas orang sebanyak 10.694 
\textit{instance} dan kelas truk sebanyak 45 \textit{instance}, atau rasio sekitar 237 
berbanding 1. Ketidakseimbangan ini bukan kekurangan proses pengumpulan data, melainkan 
cerminan komposisi objek yang sebenarnya di lapangan, sehingga dataset yang terbentuk 
dapat dikatakan merepresentasikan kondisi nyata kawasan Malioboro secara jujur.

Kedua, perbandingan delapan varian model YOLOv8 dan YOLO11 berhasil menentukan model 
dengan performa terbaik, sekaligus mengungkap pengaruh ketidakseimbangan kelas terhadap 
performa deteksi tiap kelas, menjawab rumusan masalah kedua secara menyeluruh. Pada 
himpunan uji, yolov8l mencatat mAP50 dan rerata F1 tertinggi (0,6175 dan 0,6222), 
sedangkan yolo11l mencatat mAP50-95 tertinggi (0,4012) dengan ukuran model 42\% lebih 
ringkas. Model yolo11l dipilih sebagai model terbaik berdasarkan keunggulan lokalisasi dan 
efisiensi ukuran tersebut, dan pilihan ini diperkuat oleh penyetelan resolusi citra 
masukan dari 640 menjadi 960 piksel, yang menaikkan rerata F1 model terpilih dari 0,5962 
menjadi 0,6739, melampaui capaian seluruh varian pada resolusi dasar termasuk yolov8l. 
Analisis per kelas menunjukkan ketidakseimbangan kelas berpengaruh pada performa, tetapi 
pengaruhnya dimoderasi oleh ukuran dan kekhasan visual objek, sebagaimana terlihat dari 
kelas truk yang berjumlah paling sedikit namun memperoleh F1 lebih tinggi daripada kelas 
sepeda yang berjumlah lebih banyak.

Ketiga, model terpilih berhasil diimplementasikan ke dalam sistem pemantauan kepadatan 
berbasis web yang memproses \textit{stream} video kamera CCTV dan mengukur performa 
inferensinya, menjawab rumusan masalah ketiga. Pengujian fungsional menunjukkan antarmuka 
deteksi langsung, grafik kepadatan, dan riwayat peringatan berjalan sesuai rancangan, 
dengan notifikasi yang terpicu tepat sesuai ambang kepadatan yang ditetapkan. Pengukuran 
kecepatan inferensi pada perangkat penerapan, yaitu laptop dengan GPU GTX 1650, 
menunjukkan model terpilih mencapai 9,73 \textit{frame} per detik pada resolusi 640, lebih 
cepat dibanding yolov8l yang hanya mencapai 8,80 \textit{frame} per detik pada perangkat 
yang sama, berkebalikan dengan urutan kecepatan keduanya pada GPU Tesla V100 di lingkungan 
pelatihan. Temuan ini menegaskan bahwa pertimbangan ukuran model yang mendasari pemilihan 
yolo11l pada rumusan masalah kedua terbukti relevan secara langsung bagi kecepatan pada 
perangkat target, bukan sekadar pertimbangan keamanan memori semata.

Secara keseluruhan, ketiga jawaban di atas membentuk satu rangkaian yang saling menopang. 
Dataset yang terbentuk pada jawaban pertama menjadi fondasi bagi perbandingan dan 
pemilihan model pada jawaban kedua, dan model yang terpilih serta dioptimasi pada jawaban 
kedua terbukti layak diterapkan pada sistem yang menjawab rumusan masalah ketiga, baik 
dari sisi akurasi maupun kecepatan inferensi pada perangkat target. Dengan terjawabnya 
ketiga rumusan masalah tersebut, tujuan penelitian yang ditetapkan pada Bab I tercapai 
secara menyeluruh, mulai dari ketersediaan dataset hingga sistem pemantauan yang berfungsi.

\section{Saran}
\label{sec:saran}

Penelitian ini memiliki sejumlah batasan yang ditetapkan sejak awal agar lingkupnya tetap 
terarah, sebagaimana diuraikan pada Bab I, dan beberapa keterbatasan tambahan yang 
teridentifikasi sepanjang proses penelitian. Batasan dan keterbatasan tersebut bukan 
kekurangan yang perlu disesalkan, melainkan penanda arah yang jelas bagi penelitian 
berikutnya untuk melanjutkan dari titik yang tepat. Subbab ini menguraikan sembilan arah 
pengembangan lanjutan yang diturunkan langsung dari batasan dan keterbatasan tersebut, 
disusun agar penelitian berikutnya tidak mengulang capaian yang sudah diperoleh di sini.

Validasi lapangan dan perluasan kalibrasi tingkat pelayanan. Penetapan ambang 
notifikasi berbasis tingkat pelayanan pejalan kaki pada penelitian ini bergantung pada 
estimasi luas area melalui pengukuran visual pada \textit{Google Earth}, tanpa validasi 
menggunakan pengukuran fisik langsung di lapangan, sebagaimana diakui sebagai keterbatasan 
pada Bab III. Kalibrasi ini juga hanya diterapkan pada dua titik kamera sebagai contoh 
penerapan, sedangkan titik kamera lainnya menggunakan ambang yang ditentukan secara manual. 
Penelitian lanjutan dapat melakukan pengukuran lapangan langsung, misalnya menggunakan 
meteran atau alat ukur jarak laser, untuk memvalidasi atau mengoreksi estimasi luas yang 
dihasilkan dari \textit{Google Earth}, serta mengembangkan metode estimasi luas area yang 
lebih sistematis atau bahkan otomatis untuk diterapkan pada seluruh 22 titik kamera, 
sehingga seluruh kawasan mendapat penetapan ambang yang berbasis standar kuantitatif yang 
terverifikasi, bukan hanya estimasi visual.

Penanganan kepadatan kendaraan dan kemacetan. Kesenjangan ini muncul karena penelitian 
berfokus pada kepadatan pejalan kaki melalui konsep tingkat pelayanan, sedangkan penetapan 
ambang untuk kemacetan kendaraan berada di luar lingkup penelitian, meskipun Bab I telah 
menguraikan bahwa Malioboro juga mengalami kemacetan akibat lonjakan kendaraan. Sistem 
yang dibangun sudah mendeteksi kelas kendaraan seperti mobil, motor, bus, dan truk, 
sehingga penelitian lanjutan dapat mengembangkan standar kepadatan atau tingkat pelayanan 
jalan yang setara untuk kendaraan, sehingga sistem dapat memberi notifikasi kemacetan 
selain notifikasi keramaian pejalan kaki.

Penambahan kemampuan pelacakan antar-\textit{frame}. Sistem yang dibangun melakukan 
penghitungan objek per \textit{frame} tanpa pelacakan (\textit{tracking}) antar-\textit{frame}, 
sehingga tidak dapat membedakan apakah objek yang terdeteksi pada \textit{frame} berurutan 
adalah objek yang sama atau berbeda, dan tidak dapat mengoreksi deteksi yang hilang sesaat 
akibat oklusi pada kerumunan padat sebagaimana diidentifikasi pada Bab II dan Bab IV. 
Penambahan algoritma pelacakan seperti DeepSORT pada penelitian lanjutan akan memungkinkan 
estimasi arus pergerakan, kecepatan rata-rata pejalan kaki, dan deteksi pola pergerakan 
yang tidak dapat diperoleh dari penghitungan per \textit{frame} saja, sekaligus 
menstabilkan jumlah hitungan saat objek sesaat tertutup oleh objek lain.

Verifikasi akurasi penghitungan pada kepadatan tinggi. Pengujian beban tinggi pada Bab 
IV berhasil memicu notifikasi pada ambang sesungguhnya di kamera Titik Nol dan menunjukkan 
\textit{backend} tetap stabil saat memproses lebih dari seratus objek, tetapi ketepatan 
jumlah deteksi pada kepadatan setinggi itu tidak dapat diverifikasi satu per satu karena 
kotak deteksi saling menutupi sebagian besar area \textit{frame}. Penelitian lanjutan dapat 
melakukan verifikasi akurasi penghitungan pada kondisi kepadatan tinggi secara sistematis, 
misalnya dengan membandingkan jumlah deteksi sistem terhadap hitungan manual pada potongan 
\textit{frame} terbatas atau terhadap metode \textit{crowd counting} berbasis estimasi 
kepadatan, untuk mengukur sejauh mana penghitungan tetap akurat saat oklusi antarobjek 
meningkat tajam.

Peningkatan kecepatan inferensi pada perangkat dengan sumber daya terbatas. Pengukuran 
pada Bab IV menunjukkan kecepatan inferensi model terpilih pada GPU GTX 1650 hanya 
mencapai 8,89 \textit{frame} per detik setelah penyetelan resolusi ke 960 piksel, 
turun dari 9,73 \textit{frame} per detik pada resolusi dasar, dan keduanya belum mencapai 
laju \textit{real-time} yang ideal untuk pemantauan banyak kamera secara bersamaan. 
Penelitian lanjutan dapat mengeksplorasi teknik kompresi model seperti 
\textit{pruning} atau kuantisasi, maupun konversi ke \textit{runtime} inferensi yang lebih 
efisien seperti TensorRT, untuk mempercepat inferensi pada perangkat kelas konsumen tanpa 
mengorbankan akurasi secara signifikan.

Penyetelan \textit{hyperparameter} pelatihan secara lebih luas. Penyetelan pada 
penelitian ini terbatas pada satu konfigurasi, yaitu resolusi citra masukan, sebagaimana 
diuraikan pada Bab III, sedangkan \textit{hyperparameter} pelatihan lain seperti 
\textit{learning rate}, \textit{weight decay}, dan bobot fungsi \textit{loss} tetap 
menggunakan nilai bawaan agar perbandingan antararsitektur pada Bab IV tetap adil. 
Cakupan ini memang sesuai dengan tujuan utama penelitian, tetapi belum mengeksplorasi 
apakah pencarian \textit{hyperparameter} yang lebih luas dapat meningkatkan performa model 
terpilih lebih jauh dibanding hasil penyetelan resolusi semata. Penelitian lanjutan dapat 
melakukan pencarian \textit{hyperparameter} yang lebih ekstensif pada model terpilih, 
misalnya menggunakan \textit{grid search} atau optimasi \textit{Bayesian} terhadap \textit{learning rate} 
dan bobot fungsi \textit{loss} per kelas, untuk menelaah apakah kombinasi konfigurasi 
pelatihan yang lebih luas dapat mencapai performa yang lebih tinggi dibanding 
penyetelan resolusi saja.

Penanganan ketidakseimbangan kelas pada tahap pelatihan. Bab IV menunjukkan 
ketidakseimbangan kelas memengaruhi performa deteksi pada kelas minoritas seperti sepeda 
(F1 0,270 dengan 162 \textit{instance} latih) dan bajaj (F1 0,428 dengan 96 \textit{instance} 
latih), dan penelitian ini mengatasinya secara tidak langsung melalui penyetelan resolusi 
yang menyasar objek berukuran kecil. Penelitian lanjutan dapat menambahkan strategi yang 
menyasar langsung ketidakseimbangan kelas, misalnya pembobotan fungsi kerugian per kelas 
atau augmentasi data yang lebih agresif khusus pada kelas minoritas, untuk menelaah apakah 
performa kelas seperti bajaj dan bus dapat ditingkatkan lebih lanjut tanpa mengorbankan 
performa kelas mayoritas.

Penggantian anotasi semi-otomatis dengan anotasi manual penuh. Sebagian besar data 
latih pada penelitian ini diperoleh melalui anotasi semi-otomatis, yaitu kotak pembatas 
awal dari model \textit{pretrained} YOLOv8x dan YOLO11x yang dikoreksi secara manual, 
sebagaimana diuraikan pada Bab III. Pendekatan ini mempercepat proses anotasi, tetapi kotak 
pembatas awal yang dihasilkan tetap mengikuti karakteristik visual dataset asal model 
\textit{pretrained} tersebut, sehingga ada kemungkinan kotak yang dianggap benar oleh 
anotator pada tahap koreksi sebenarnya kurang merepresentasikan kondisi objek yang 
sesungguhnya di lapangan, khususnya untuk kelas yang bentuknya berbeda dari kelas serupa 
pada dataset umum seperti becak. Penelitian lanjutan dapat membangun ulang dataset 
menggunakan anotasi manual penuh tanpa bantuan model \textit{pretrained}, sehingga setiap 
kotak pembatas benar-benar mengikuti kondisi visual objek di kawasan Malioboro, dan menguji 
apakah performa model yang dilatih pada dataset tersebut lebih tinggi dibanding performa 
yang dicapai pada dataset hasil anotasi semi-otomatis.

Pengujian keterlibatan pengguna dan kebergunaan sistem. Penelitian ini berfokus pada 
perbandingan model deteksi dan implementasi fungsional sistem, sedangkan perancangan 
berbasis \textit{User-Centered Design} dan pengujian \textit{usability} tidak termasuk 
dalam lingkup penelitian. Penelitian lanjutan dapat melibatkan calon pengguna sesungguhnya, 
seperti petugas pengelola kawasan Malioboro, untuk menilai kemudahan penggunaan antarmuka 
dan menyesuaikan rancangan sistem berdasarkan masukan tersebut, sehingga sistem tidak 
hanya terbukti berfungsi secara teknis tetapi juga nyaman digunakan dalam operasional 
sehari-hari.
